\documentclass[11pt]{article}

\usepackage[margin=1in]{geometry}
\usepackage{amsmath,amssymb,amsthm,mathtools}
\usepackage[hidelinks]{hyperref}

\newtheorem{theorem}{Theorem}[section]
\newtheorem{proposition}[theorem]{Proposition}
\newtheorem{corollary}[theorem]{Corollary}
\newtheorem{remark}[theorem]{Remark}

\newcommand{\Rn}{\mathbb{R}^{n}}
\newcommand{\BR}{B_R}
\newcommand{\norm}[1]{\left\lVert #1\right\rVert}

\title{E68: Localized Gradient-Energy Estimate\\
\large Counterexamples, Dimensional Classification, and Corrected Theorems}
\author{Mathematical audit of WCT Manuscript 4}
\date{October 26, 2026}

\begin{document}
\maketitle

\begin{abstract}
The proposed E68 estimate bounds the gradient energy inside a ball by
$CR^{n-2}\norm{\psi}_{H^1(\Rn)}^2$. We prove that, with a constant
independent of radius and field, this assertion holds for all $H^1$ fields
and all radii exactly when $n=2$. For every $n>2$, a single compactly
supported, singular $H^1$ field violates the estimate at arbitrarily small
radii; a family of smooth compactly supported fields also rules out a
universal constant. In dimension $n=1$ the proposed estimate fails at
large radii. We then prove a corrected inequality under the assumption
$\nabla\psi\in L^n$ and obtain its $H^2$ consequence in dimensions three
and four. These results concern analytic energy estimates, not the
existence or dynamical stability of confined wave solutions.
\end{abstract}

\section{Statement and quantifiers}

For $n\geq1$, let $B_R(x_0)=\{x\in\Rn:|x-x_0|<R\}$, and use the
inhomogeneous norm
\[
\norm{\psi}_{H^1(\Rn)}^2
=\norm{\psi}_{L^2(\Rn)}^2+\norm{\nabla\psi}_{L^2(\Rn)}^2.
\]
The E68 assertion, interpreted as a uniform estimate, is
\begin{equation}\label{eq:E68}
\boxed{\qquad
\int_{B_R(0)}|\nabla\psi(x)|^2\,dx
\leq C_n R^{n-2}\norm{\psi}_{H^1(\Rn)}^2,
\qquad \psi\in H^1(\Rn),\ R>0.
\qquad}
\tag{E68}
\end{equation}
Here $C_n<\infty$ must be independent of $R$ and $\psi$. The same
conclusions hold if one requires the estimate uniformly over all centers
$x_0\in\Rn$.

\begin{theorem}[Exact dimensional classification]\label{thm:classification}
For the class $H^1(\Rn)$ and all radii $R>0$, a uniform constant in
\eqref{eq:E68} exists if and only if $n=2$. In that dimension $C_2=1$
works. For each $n>2$, there is even a single
$\psi\in H^1(\Rn)$ for which no finite constant depending on $\psi$
can make \eqref{eq:E68} hold for all sufficiently small $R$.
\end{theorem}

\begin{proof}
We consider the three cases.

\medskip\noindent
\textbf{Case $n=2$.} Since $R^{n-2}=1$ and $B_R\subset\mathbb R^2$,
\[
\int_{B_R}|\nabla\psi|^2\,dx
\leq\int_{\mathbb R^2}|\nabla\psi|^2\,dx
\leq\norm{\psi}_{H^1(\mathbb R^2)}^2.
\]
Thus the inequality holds with $C_2=1$.

\medskip\noindent
\textbf{Case $n=1$.} Choose a nonconstant
$\psi\in C_c^\infty((-1,1))$ and put
$A=\int_{\mathbb R}|\psi'|^2\,dx>0$.
For every $R\geq1$,
\[
\int_{-R}^{R}|\psi'|^2\,dx=A.
\]
If \eqref{eq:E68} were true, then
\[
A\leq C_1R^{-1}\norm{\psi}_{H^1(\mathbb R)}^2
\qquad(R\geq1).
\]
Sending $R\to\infty$ makes the right-hand side tend to zero, a
contradiction.

\medskip\noindent
\textbf{Case $n>2$.} Fix
\[
0<a<\frac{n-2}{2},
\]
and choose $\chi\in C_c^\infty([0,\infty))$ with
$\chi(r)=1$ for $0\leq r\leq1$ and $\chi(r)=0$ for $r\geq2$.
Define, for $x\neq0$,
\[
\psi(x)=\chi(|x|)|x|^{-a},
\]
with an arbitrary value at $x=0$. Write
$\sigma_{n-1}=|\mathbb S^{n-1}|$. Since $\psi=r^{-a}$ in $B_1$,
polar coordinates yield
\begin{equation}\label{eq:L2integral}
\int_{B_1}|\psi|^2\,dx
=\sigma_{n-1}\int_0^1r^{n-1-2a}\,dr
=\frac{\sigma_{n-1}}{n-2a}<\infty.
\end{equation}
For $0<|x|<1$,
\[
\nabla\psi(x)=-a|x|^{-a-2}x,
\qquad |\nabla\psi(x)|^2=a^2|x|^{-2a-2},
\]
and therefore
\begin{equation}\label{eq:gradintegral}
\int_{B_1}|\nabla\psi|^2\,dx
=\sigma_{n-1}a^2\int_0^1r^{n-3-2a}\,dr
=\frac{\sigma_{n-1}a^2}{n-2-2a}<\infty.
\end{equation}
The cutoff-annulus contributions are finite, and $\psi$ vanishes
outside $B_2$. These classical derivatives on $\Rn\setminus\{0\}$
are also the weak derivatives: applying integration by parts outside
$B_\varepsilon$ leaves a boundary term bounded by a constant times
$\varepsilon^{n-1-a}$, which tends to zero because
$a<(n-2)/2<n-1$. Thus $\psi\in H^1(\Rn)$.

For each $0<R<1$, the calculation in \eqref{eq:gradintegral} gives
\begin{equation}\label{eq:localexact}
\int_{B_R}|\nabla\psi|^2\,dx
=\frac{\sigma_{n-1}a^2}{n-2-2a}R^{n-2-2a}.
\end{equation}
Hence
\begin{equation}\label{eq:blowup}
R^{2-n}\int_{B_R}|\nabla\psi|^2\,dx
=\frac{\sigma_{n-1}a^2}{n-2-2a}R^{-2a}
\longrightarrow+\infty
\qquad(R\downarrow0).
\end{equation}
On the other hand, the proposed E68 bound would make the left-hand
side at most $C_n\norm{\psi}_{H^1}^2<\infty$, independently of $R$.
This contradiction completes the proof.
\end{proof}

\section{Failure on smooth fields}

The preceding fixed-field counterexample has a singularity at the
origin. The failure of a universal constant does not require any
singular fields.

\begin{proposition}[Smooth concentration counterexample]
\label{prop:smooth}
Let $n>2$. For every $C>0$, there are
$\psi\in C_c^\infty(\Rn)$ and $R>0$ for which
\[
\int_{B_R}|\nabla\psi|^2\,dx
>C R^{n-2}\norm{\psi}_{H^1(\Rn)}^2.
\]
\end{proposition}

\begin{proof}
Choose a nonconstant $\phi\in C_c^\infty(B_1)$ and put
\[
\psi_R(x)=R^{1-n/2}\phi(x/R).
\]
Let
\[
A=\norm{\nabla\phi}_2^2>0,
\qquad B=\norm{\phi}_2^2.
\]
Differentiation and the substitution $y=x/R$ give
\[
\norm{\nabla\psi_R}_2^2=A,
\qquad \norm{\psi_R}_2^2=R^2B.
\]
Since $\operatorname{supp}\psi_R\subset B_R$,
\[
\int_{B_R}|\nabla\psi_R|^2\,dx=A,
\qquad \norm{\psi_R}_{H^1}^2=A+R^2B.
\]
A uniform E68 inequality would imply
\[
A\leq CR^{n-2}(A+R^2B)
\qquad\text{for every }R>0.
\]
The right-hand side tends to zero as $R\downarrow0$, since $n>2$,
while $A>0$ is fixed. This is impossible.
\end{proof}

\section{Corrected estimate with critical gradient integrability}

\begin{theorem}[Corrected E68]\label{thm:corrected}
Let $n>2$ and suppose $\psi\in H^1(\Rn)$ with
$\nabla\psi\in L^n(\Rn)$. For every $R>0$ and $x_0\in\Rn$,
\begin{equation}\label{eq:corrected}
\boxed{\quad
\int_{B_R(x_0)}|\nabla\psi|^2\,dx
\leq\omega_n^{1-2/n}R^{n-2}
\norm{\nabla\psi}_{L^n(\Rn)}^2,
\quad}
\tag{E68$'$}
\end{equation}
where $\omega_n=|B_1(0)|$.
\end{theorem}

\begin{proof}
Apply H\"older's inequality on $B_R(x_0)$ with conjugate
exponents $p=n/2$ and $q=n/(n-2)$:
\begin{align*}
\int_{B_R(x_0)}|\nabla\psi|^2\,dx
&=\int_{B_R(x_0)}|\nabla\psi|^2\cdot1\,dx\\
&\leq
\left(\int_{B_R(x_0)}|\nabla\psi|^{2p}\,dx\right)^{1/p}
\left(\int_{B_R(x_0)}1^q\,dx\right)^{1/q}\\
&=\left(\int_{B_R(x_0)}|\nabla\psi|^n\,dx\right)^{2/n}
|B_R(x_0)|^{(n-2)/n}\\
&\leq\norm{\nabla\psi}_{L^n(\Rn)}^2
(\omega_nR^n)^{(n-2)/n}\\
&=\omega_n^{1-2/n}R^{n-2}
\norm{\nabla\psi}_{L^n(\Rn)}^2.
\end{align*}
This is \eqref{eq:corrected}.
\end{proof}

\begin{corollary}[Recovering the original $H^1$ right-hand side]
\label{cor:apriori}
Let $n>2$ and let $\mathcal A\subset H^1(\Rn)$ be an admissible
class satisfying the \emph{independently established} a priori bound
\[
\norm{\nabla\psi}_{L^n(\Rn)}
\leq K\norm{\psi}_{H^1(\Rn)}
\qquad\text{for all }\psi\in\mathcal A,
\]
with $K$ independent of $\psi$. Then
\[
\int_{B_R(x_0)}|\nabla\psi|^2\,dx
\leq\omega_n^{1-2/n}K^2R^{n-2}
\norm{\psi}_{H^1(\Rn)}^2
\qquad (\psi\in\mathcal A).
\]
\end{corollary}

\begin{proof}
Insert the assumed gradient estimate into
Theorem~\ref{thm:corrected}.
\end{proof}

\begin{remark}
The hypothesis in Corollary~\ref{cor:apriori} does not follow from
$H^1$ membership alone. A governing PDE, a band limit, or an additional
regularity restriction may supply such a bound, but it must be proved
separately for the specified admissible class.
\end{remark}

\section{An $H^2$ replacement in dimensions three and four}

\begin{corollary}[Higher-regularity localized estimate]
\label{cor:H2}
If $n\in\{3,4\}$ and $\psi\in H^2(\Rn)$, there is a constant
$C_n<\infty$, independent of $\psi$, $x_0$, and $R$, such that
\begin{equation}\label{eq:H2}
\boxed{\quad
\int_{B_R(x_0)}|\nabla\psi|^2\,dx
\leq C_nR^{n-2}\norm{\psi}_{H^2(\Rn)}^2.
\quad}
\end{equation}
\end{corollary}

\begin{proof}
When $n=4$, the standard Sobolev inequality in $\mathbb R^4$
applied to $\nabla\psi\in H^1(\mathbb R^4)$ gives
\[
\norm{\nabla\psi}_{L^4}
\leq C\norm{\nabla\psi}_{H^1}
\leq C\norm{\psi}_{H^2}.
\]
When $n=3$, Sobolev embedding yields
$\norm{\nabla\psi}_{L^6}\leq C\norm{\psi}_{H^2}$.
Interpolation between $L^2$ and $L^6$ gives
\[
\norm{\nabla\psi}_{L^3}
\leq\norm{\nabla\psi}_{L^2}^{1/2}
\norm{\nabla\psi}_{L^6}^{1/2}
\leq C\norm{\psi}_{H^2}.
\]
In both dimensions the hypothesis of Theorem~\ref{thm:corrected}
holds with $\norm{\nabla\psi}_{L^n}\leq C\norm{\psi}_{H^2}$.
Its conclusion is precisely \eqref{eq:H2}.
\end{proof}

\begin{proposition}[The $H^2$ replacement fails above dimension four]
\label{prop:H2fail}
If $n>4$, there is no constant $C_n$ for which
\eqref{eq:H2} holds for all $\psi\in C_c^\infty(\Rn)$ and all $R>0$.
\end{proposition}

\begin{proof}
Take a nonconstant $\phi\in C_c^\infty(B_1)$ and set
\[
u_R(x)=R^{2-n/2}\phi(x/R),\qquad 0<R<1.
\]
Scaling gives
\[
\norm{u_R}_2^2=R^4\norm{\phi}_2^2,
\quad \norm{\nabla u_R}_2^2=R^2\norm{\nabla\phi}_2^2,
\quad \norm{D^2u_R}_2^2=\norm{D^2\phi}_2^2.
\]
Thus $\norm{u_R}_{H^2}^2$ stays bounded as $R\downarrow0$, whereas
\[
\int_{B_R}|\nabla u_R|^2\,dx
=R^2\norm{\nabla\phi}_2^2.
\]
If \eqref{eq:H2} held universally, it would imply
\[
R^2\norm{\nabla\phi}_2^2\leq C R^{n-2}
\]
for some fixed $C$ and every small $R$. Dividing by $R^2$ gives
$\norm{\nabla\phi}_2^2\leq C R^{n-4}\to0$, a contradiction.
\end{proof}

\section{Status and scope for Manuscript 4}

\begin{itemize}
\item \textbf{REFUTED:} E68 with its unrestricted $H^1$ hypothesis,
when $n>2$, even if the multiplicative constant is allowed to depend on
one fixed singular $H^1$ field. The all-radii $n=1$ version also fails.
\item \textbf{PROVED:} E68 in dimension $n=2$ with $C_2=1$.
\item \textbf{PROVED:} Corrected E68$'$ when $n>2$ and
$\nabla\psi\in L^n$, with the displayed explicit constant.
\item \textbf{PROVED:} The $H^2$-based estimate for $n=3,4$.
\item \textbf{CONDITIONAL:} Recovery of the original $H^1$ right-hand
side within a PDE-dependent admissible class, pending an independent
uniform $L^n$ gradient estimate.
\end{itemize}

These are purely analytic assertions. They do not establish existence,
orbital stability, metastability, or a physical confinement-lifetime
comparison in different spatial dimensions.

\begin{thebibliography}{9}
\bibitem{WCTCorrected}
\emph{WCT Full Equation List (Corrected)}, E68 entry,
\url{https://github.com/rickyjreyes/geometry_of_resonance/blob/main/WCT_FULL_EQUATION_LIST_CORRECTED.md}.

\bibitem{WCTLean}
\emph{WCT Lean Analytic Contracts}, definition \texttt{LocalEnergyBound},
\url{https://github.com/rickyjreyes/wct-lean/blob/main/WCTLean/Contracts/Analytic.lean}.
\end{thebibliography}

\end{document}
